\section{The Response That Is Not There}
\label{sec:ch14-the-response-that-is-not-there}

\index{blast radius!the whole response}%
The Ratings service has a root field of its own.
Chapter~\ref{ch:second-third-subgraph} gave it one because every subgraph needs
somewhere to put a query, and almost everything interesting that service knows
is about a type it does not own. What was left was a total:

\begin{graphqlsdl}
type Query {
  "How many scores this service holds, across every session."
  ratingCount: Int!
}
\end{graphqlsdl}

\index{non-null!a count that cannot fail}%
That exclamation mark is defensible on every rule this book has given so far. A
\code{COUNT} over a table this service owns is a value it can always produce.
There is no missing row to worry about, no foreign key that might not match, no
other team involved. Chapter~\ref{ch:schema-design} said to mark a field
non-null when you can guarantee it from data you own, in the same process, with
a constraint behind it, and this field clears that bar without an argument.

Stop the service and ask for it.

\begin{minted}{text}
{"query":"{ ratingCount }"}
\end{minted}

\begin{minted}[fontsize=\footnotesize,breakanywhere]{json}
{"errors":[{"message":"Failed to fetch from Subgraph 'ratings'."}],"data":null}
\end{minted}

\index{null propagation!reaching the data entry}%
\code{data} is null. Not \code{data.ratingCount}: the entry itself. There was
nowhere else for it to go, because the parent of a root field is the response.

\subsection{And It Takes the Rest With It}

\index{blast radius!another service's answer}%
That is a small loss while the request asks for one field. Ask for two:

\begin{minted}{text}
{"query":"{ ratingCount sessions(first: 10) { nodes { title } } }"}
\end{minted}

\begin{minted}[fontsize=\footnotesize,breakanywhere]{json}
{"errors":[{"message":"Failed to fetch from Subgraph 'ratings'."},
{"message":"Cannot return null for non-nullable field 'Query.ratingCount'.",
"path":["ratingCount"]}],
"data":null}
\end{minted}

\index{statement count!the healthy service still paid}%
The Sessions service is running. It was asked for the schedule, it issued its
one statement, it read four rows, and every title was correct. None of that is
in the response. A field on a stopped service, in a different process, owned by
a different team, on a different database, emptied it.

\index{blast radius!the two positions that matter}%
In section~\ref{sec:ch14-four-sessions-for-one-dead-service} a non-null field
inside somebody else's list cost that list. Here a non-null field at the root
costs the response. Those are the only two sizes of collateral damage a
federated graph has, and which one you get is decided by where in the
supergraph the field sits rather than by anything about the failure.

\index{HTTP 200!with no data at all}%
Two hundred, incidentally, as it is for every response in this chapter. The
router refuses a request outright only when it cannot verify a token, which is
the 401 chapter~\ref{ch:entities-authorization} met; a graph that cannot answer
is a success as far as the transport is concerned. A client that checks the
status code sees one here, with no data in it.

\subsection{What This Costs a Real Graph}

\index{blast radius!at production scale}%
Four sessions and a rating count is a small demonstration of a large problem.
The graph in front of you has four services. A graph with forty has forty
chances that somebody, on some team, marked a root field non-null for a reason
that was locally correct, and any one of them turns a partial outage into a
total one for every client whose query happens to mention that field.

\index{ownership!a promise made in another team's file}%
Nobody reviewing the Ratings service would catch it either. The field is a
count over its own table. The reviewer who approves it is right about
everything in front of them, and the consequence is in a file they have never
opened, in a response shape that does not exist until a composer builds it.
